\scanpage{10}\mspage{9}

\uline{NB} Pour bien faire, il faudrait prouver que sur le site modulaire $S$ des schémas en courbes \inserted{relatives propres lisses connexes} de genre $g$\deleted{$\unsure{(g\geq1)|\geq1}$} \inserted{($g\geq1$ ; on \deleted{\unsure{distingue}}\inserted{\unsure{fixe|prend|choisit}} une section marquée pour $g=1$)} \deleted{\unsure{le|tel que}} si $\mathcal X$ est la courbe de genre $g\deleted{\unsure{*}}$ universelle, $P=\underline{\mathrm{Pic}}_{X/S}$, alors le groupe des sections de $P$ sur $S$ est engendré par $K$ (\deleted{donc} \inserted{en particulier, que} le groupe des sections de $P^{\unsure{\circ|0}}$ sur $S$ est égal à zéro). Il suffirait de déterminer le groupe en question en car. nulle.

\deleted{\unsure{f)}} \inserted{f) \uline{Ex. 2}} \uline{Cas} $X$ \uline{un schéma abélien}. Donc $X=A$. On sait que \inserted{sur fibres géométriques} $\exp D$ définit $\deleted{\unsure{A\to A}}\inserted{\Phi}:\mathrm{Chow}(A)_{\mathbb Q}\xrightarrow{\sim}\mathrm{Chow}(B)_{\mathbb Q}$ qui \deleted{\unsure{en}} transforme élément de degré $i$ en élément de degré complémentaire $n-i$, \inserted{et} qui transforme produit ordinaire en produit de convolution, et inversement. D'ailleurs $T_{X/k}=1$, donc on trouve

\[
\varphi(\delta)=\inserted{\text{composante de degré 1 de }}\sum\tfrac1{i!}\underbrace{\Phi(\Delta)}_{\substack{\text{classe d'une}\\\text{courbe}}}^{*i}
=\frac1{(n-1)!}\Phi(\Delta)^{*(n-1)}.
\]

\deleted{\unsure{donc}} On trouve ainsi un hom (\inserted{polynomial} de degré $n-1$)

\[
\varphi:\underline{\mathrm{NS}}_{\deleted{\unsure{X/S}}A/S}\longrightarrow\underline{\mathrm{NS}}_{B/S}.
\]

On vérifie que, si $\delta$ est une section de $\underline{\mathrm{NS}}_{A/S}$, définissant $\widetilde\delta:A\longrightarrow B$ d'où $\varphi(\delta)$ section de $\underline{\mathrm{NS}}_{\unsure{B|P}/S}$, d'où

\[
\widetilde{\varphi(\delta)}:B\longrightarrow A
\]

\margin{faire calcul direct, en faisant \deleted{\unsure{attention}} attention aux signes !}

\[
\text{on a}\qquad\boxed{\widetilde{\varphi(\delta)}\,\widetilde\delta=-\chi(\delta)\,\mathrm{id}_A,\qquad\widetilde\delta\,\widetilde{\varphi(\delta)}=-\chi(\delta)\,\mathrm{id}_B}
\]
